% !TEX root = ../main.tex
\section{Optical Transitions: Absorption, Stimulated and Spontaneous Emission}
\label{sec:optical_transitions}

\begin{figure}[H]
	\centering
	\begin{tikzpicture}[
		x=1cm,
		y=1cm,
		scale=0.95,
		every node/.style={font=\small\setstretch{1}},
		band/.style={line width=1.6pt},
		level/.style={dashed},
		transition/.style={-{Stealth[length=2.5mm]}, semithick},
		photon/.style={-{Stealth[length=2.5mm]}, semithick, decorate, decoration={snake, amplitude=1.1mm, segment length=3mm, pre length=2mm, post length=2.5mm}},
		electron/.style={circle, fill, inner sep=2.2pt},
		hole/.style={circle, draw, semithick, fill=white, inner sep=2.2pt},
	]
		% spontaneous emission
		\draw[level] (0.1, 2.4) -- (1.3, 2.4);
		\draw[level] (0.1, 0) -- (1.3, 0);
		\draw[{Stealth[length=2.5mm]}-{Stealth[length=2.5mm]}, semithick]
			(0.5, 0) --
			node[right] {$E_{\mathrm{g}}$}
			(0.5, 2.4);
		\draw[band] (1.9, 2.4) -- (4.3, 2.4);
		\draw[band] (1.9, 0) -- (4.3, 0);
		\draw[transition] (2.3, 3.0) -- (2.3, -0.62);
		\node[electron] at (2.3, 3.0) {};
		\node[hole] at (2.3, -0.8) {};
		\draw[photon] (2.6, 1.2) -- node[above=1.5mm] {Photon} (4.3, 1.2);
		\node[align=center] at (0.7, 3.25) {Conduction\\band};
		\node[align=center] at (0.7, -0.95) {Valence\\band};
		\node[anchor=base] at (3.1, -2.0) {spontaneous emission};

		% absorption
		\draw[band] (5.5, 2.4) -- (7.9, 2.4);
		\draw[band] (5.5, 0) -- (7.9, 0);
		\draw[transition] (7.5, -0.62) -- (7.5, 2.98);
		\node[electron] at (7.5, 3.0) {};
		\node[hole] at (7.5, -0.8) {};
		\draw[photon] (5.5, 1.2) -- (7.2, 1.2);
		\node[left=3mm] at (7.5, 3.0) {Electron};
		\node[left=3mm] at (7.5, -0.8) {Hole};
		\node[anchor=base] at (6.7, -2.0) {absorption};

		% stimulated emission
		\draw[band] (9.1, 2.4) -- (13.3, 2.4);
		\draw[band] (9.1, 0) -- (13.3, 0);
		\draw[transition] (11.2, 3.0) -- (11.2, -0.62);
		\node[electron] at (11.2, 3.0) {};
		\node[hole] at (11.2, -0.8) {};
		\draw[photon] (9.1, 1.2) -- (10.9, 1.2);
		\draw[photon] (11.5, 1.65) -- (13.3, 1.65);
		\draw[photon] (11.5, 0.75) -- (13.3, 0.75);
		\node[anchor=base] at (11.2, -2.0) {stimulated emission};

		% transition energy
		\draw[level] (13.6, 3.0) -- (14.8, 3.0) node[above left] {$E_{2}$};
		\draw[level] (13.6, -0.8) -- (14.8, -0.8) node[below left] {$E_{1}$};
		\draw[{Stealth[length=2.5mm]}-{Stealth[length=2.5mm]}, semithick]
			(14.0, -0.8) --
			node[right] {$\hbar\omega$}
			(14.0, 3.0);
	\end{tikzpicture}
	\caption{Band diagrams of the optical transitions between the conduction and
	the valence band: spontaneous emission, absorption and stimulated emission. \imcite{benndorf_semiconductor_nodate}}
	\label{fig:optical_transitions}
\end{figure}

An incoming photon induces a perturbation of the electronic system of the
semiconductor, which may result in the generation of an electron-hole pair. This
process is called absorption. For this electronic excitation to occur, the
energy $\hbar \omega$ of the incoming photon must be at least the band gap
energy $E_{\mathrm{g}}$.

An incoming photon induces a perturbation of the electronic system of the
semiconductor, which may result in a radiative recombination of an electron-hole pair.
The incoming and the generated photon are coherent, i.e., they share the same direction,
phase, frequency and polarization. The process requires an electron in the
conduction band and a hole in the valence band and is called stimulated emission.

Without any incoming photon, an electron in the conduction band can recombine radiatively
with a hole in the valence band. The emitted photons have random directions and
phases, so the resulting light is incoherent. This process is called spontaneous
emission. All three processes are depicted in \cref{fig:optical_transitions}.

\subsection{Mathematical Description}
\label{sec:mathematical_description}

\begin{table}
	\centering
	\begin{tabular}{lll}
		\toprule Process            & $i$             & $\Phi^{(i)}$                           \\
		\midrule absorption (gross) & $\mathrm{abs}$  & $(1-f_{\mathrm{n}})(1-f_{\mathrm{p}})$ \\
		stimulated emission (gross) & $\mathrm{stim}$ & $f_{\mathrm{n}}f_{\mathrm{p}}$         \\
		spontaneous emission        & $\mathrm{spon}$ & $f_{\mathrm{n}}f_{\mathrm{p}}$         \\
		\bottomrule
	\end{tabular}
	\caption{Occupation factors for absorption, stimulated and spontaneous emission.}
	\label{tab:occupation_factors}
\end{table}

Absorption, stimulated emission and spontaneous emission can all be described
using a single formula based on Fermi's Golden Rule. The rate coefficients $r^{(i)}$,
i.e., the rates normalized to one photon per mode, are given by:
\begin{equation}
	r^{(i)}(\hbar\omega)=P(\hbar\omega)\sum_{\mathbf{k}_{\mathrm{p}},\mathbf{k}_{\mathrm{n}}}
	\frac{2}{V}\lvert\hat e\cdot\mathbf{p}_{\mathrm{cv}}(\mathbf{k}_{\mathrm{p}},\mathbf{k}
	_{\mathrm{n}})\rvert^{2}\,\delta\big(E_{\mathrm{c}}(\mathbf{k}_{\mathrm{n}})-E_{\mathrm{v}}
	(\mathbf{k}_{\mathrm{p}})-\hbar\omega\big)\;\Phi^{(i)}\label{eq:universal_rate_equation}
\end{equation}
The factor $2/V$ counts both spin orientations per unit volume. $\Phi^{(i)}$ is the
occupation factor of process $i$, eludicated in \cref{tab:occupation_factors}. The
occupation factors depend on $E_{\mathrm{c}}(\mathbf{k}_{\mathrm{n}})$ and $E_{\mathrm{v}}
(\mathbf{k}_{\mathrm{p}})$ individually (via $f_{\mathrm{n}}$ and
$f_{\mathrm{p}}$), not only on their difference. The remaining prefactors are given
by:
\begin{equation}
	P(\hbar\omega)=\frac{n_{\mathrm{r}}^{2}(\hbar\omega)^{2}}{\pi^{2}\hbar^{3}c^{2}}
	\,C_{0}(\hbar\omega),\qquad C_{0}(\hbar\omega)=\frac{\pi e^{2}\hbar}{n_{\mathrm{r}}c\,\epsilon_{0}m_{0}^{2}\,\hbar\omega}
	\quad\Rightarrow\quad P(\hbar\omega)=\frac{n_{\mathrm{r}}e^{2}}{\pi\hbar^{2}c^{3}\epsilon_{0}m_{0}^{2}}
	\;\hbar\omega.
\end{equation}

Note that $r^{(i)}$ are normalized to one photon per mode and can therefore be regarded
as material-specific quantities. However, the actual absorption and stimulated emission
rates depend on the photon occupation number $n_{\mathrm{ph}}$:
\begin{equation}
	R^{\mathrm{abs}}= n_{\mathrm{ph}}\cdot r^{\mathrm{abs}}, \qquad R^{\mathrm{stim}}
	= n_{\mathrm{ph}}\cdot r^{\mathrm{stim}}, \qquad R^{\mathrm{spon}}=r^{\mathrm{spon}}
	.
\end{equation}

\subsection{Analytical Line Shape}
\label{sec:line_shape}

Under three additional assumptions, the universal emission and absorption rates yield
an analytical line shape:

\begin{enumerate}[label=(\roman{*}), ref=(\roman{*})]
	\item \label{ass:k_selection} Vertical transitions ($\mathbf{k}$-selection):
		$\lvert \hat{e}\cdot \mathbf{p}_{\mathrm{cv}}(\mathbf{k}_{\mathrm{p}},\mathbf{k}
		_{\mathrm{n}}) \rvert^{2}=\lvert \hat{e}\cdot \mathbf{p}_{\mathrm{cv}}\rvert^{2}
		\, \delta_{\mathbf{k}_{\mathrm{p}} \mathbf{k}_{\mathrm{n}}}$,

	\item \label{ass:const_matrix} a constant matrix element:
		$\lvert\hat e\cdot\mathbf{p}_{\mathrm{cv}}\rvert^{2}=\mathrm{const.}$,

	\item \label{ass:parabolic} parabolic bands, see \cref{eq:parabolic_bands}.
\end{enumerate}

Substituting assumptions~\ref{ass:k_selection} and~\ref{ass:const_matrix} into
\cref{eq:universal_rate_equation} yields:
\begin{equation}
	r^{(i)}(\hbar\omega)= P(\hbar\omega)\, \frac{2}{V}\lvert\hat e\cdot\mathbf{p}_{\mathrm{cv}}
	\rvert^{2}\sum_{\mathbf{k}}\delta\big(E_{\mathrm{cv}}(\mathbf{k})-\hbar\omega\big
	)\;\Phi^{(i)}(\mathbf{k})
\end{equation}

The sum is converted to an integral with the density of $\mathbf{k}$-states
$V/(2 \pi)^{3}$.
\begin{equation}
	\quad r^{(i)}(\hbar\omega)= P(\hbar\omega)\,\lvert\hat e\cdot\mathbf{p}_{\mathrm{cv}}
	\rvert^{2}\, \frac{2}{(2\pi)^{3}}\int \delta\big(E_{\mathrm{cv}}(\mathbf{k})-\hbar
	\omega\big)\;\Phi^{(i)}(\mathbf{k})\,\mathrm{d}\mathbf{k}
\end{equation}

For the isotropic bands of assumption~\ref{ass:parabolic}, the integrand depends only
on $k=\lvert\mathbf{k}\rvert$, so the angular integration can be performed
directly and $\mathrm{d}\mathbf{k}\to4\pi k^{2}\,\mathrm{d}k$:
\begin{equation}
	r^{(i)}(\hbar\omega)= P(\hbar\omega)\,\lvert\hat e\cdot\mathbf{p}_{\mathrm{cv}}
	\rvert^{2}\, \frac{1}{\pi^{2}}\int_{0}^{\infty}k^{2}\,\delta\big(E_{\mathrm{cv}}
	(k)-\hbar\omega\big)\;\Phi^{(i)}(k)\,\mathrm{d}k
\end{equation}

Changing the integration variable from $k$ to $E_{\mathrm{cv}}$ with the
parabolic dispersion of assumption~\ref{ass:parabolic} yields: With $k=k(E_{\mathrm{cv}}
)$, the occupation factor becomes $\Phi^{(i)}\big(k(E_{\mathrm{cv}})\big)$, and the
delta distribution can be evaluated directly. It fixes
$E_{\mathrm{cv}}(k_{0})=\hbar\omega$, i.e., $k=k_{0}$ with $k_{0}$ the magnitude of
the wavevector at photon energy $\hbar\omega$, so that $\Phi^{(i)}$ is evaluated
at $k_{0}$:
\begin{align}
	r^{(i)}(\hbar\omega) & = P(\hbar\omega)\,\lvert\hat e\cdot\mathbf{p}_{\mathrm{cv}}\rvert^{2}\, \frac{1}{2\pi^{2}}\left(\frac{2m_{\mathrm{r}}}{\hbar^{2}}\right)^{3/2}\int_{E_{\mathrm{g}}}^{\infty}\sqrt{E_{\mathrm{cv}}-E_{\mathrm{g}}}\;\delta(E_{\mathrm{cv}}-\hbar\omega)\;\Phi^{(i)}\big(k(E_{\mathrm{cv}})\big)\,\mathrm{d}E_{\mathrm{cv}} \\
	                     & = P(\hbar\omega)\,\lvert\hat e\cdot\mathbf{p}_{\mathrm{cv}}\rvert^{2}\, \underbrace{\frac{1}{2\pi^{2}}\left(\frac{2m_{\mathrm{r}}}{\hbar^{2}}\right)^{3/2}\sqrt{\hbar\omega-E_{\mathrm{g}}}\;\Theta(\hbar\omega-E_{\mathrm{g}})}_{D_{\mathrm{j}}(\hbar\omega)}\;\Phi^{(i)}(k_{0})
\end{align}

In this equation, $D_{\mathrm{j}}$ is the Joint Density of States (including
spin). With the energy zero at the valence-band maximum, $E_{\mathrm{c}}(k)=E_{\mathrm{g}}
+\frac{\hbar^{2}k^{2}}{2m_{\mathrm{e}}^{*}}$ and $E_{\mathrm{v}}(k)=-\frac{\hbar^{2}k^{2}}{2m_{\mathrm{h}}^{*}}$,
the occupation factors are evaluated at $k_{0}$, i.e., at the band energies $E_{\mathrm{c}}
(k_{0})$ and $E_{\mathrm{v}}(k_{0})$:
\begin{equation}
	k_{0}=\frac{\sqrt{2m_{\mathrm{r}}(\hbar\omega-E_{\mathrm{g}})}}{\hbar},\qquad E
	_{\mathrm{c}}(k_{0})=E_{\mathrm{g}}+\frac{m_{\mathrm{r}}}{m_{\mathrm{e}}^{*}}(\hbar
	\omega-E_{\mathrm{g}}),\qquad E_{\mathrm{v}}(k_{0})=-\frac{m_{\mathrm{r}}}{m_{\mathrm{h}}^{*}}
	(\hbar\omega-E_{\mathrm{g}})
\end{equation}

The absorption and emission line shape therefore starts at $E_{\mathrm{g}}$ and rises
as $\sqrt{\hbar\omega-E_{\mathrm{g}}}$, modulated by the prefactor
$P(\hbar \omega)$ and the occupation factor $\Phi^{(i)}(k_{0})$.

\subsection{Rate Ratios}
\label{sec:rate_ratios}

As absorption, stimulated emission and spontaneous emission follow from the same
general \cref{eq:universal_rate_equation}, their ratios can be calculated easily.

\subsubsection{Stimulated Emission vs Absorption}

Let $f_{\mathrm{n}}$ and $f_{\mathrm{p}}$ be the occupation probabilities of the
conduction band by electrons and of the valence band by holes, described by the Quasi
Fermi Levels $F_{\mathrm{n}}$ and $F_{\mathrm{p}}$ at the same temperature $T$:
\begin{align}
	\frac{1-f_{\mathrm{n}}}{f_{\mathrm{n}}}=\exp\left(\frac{E_{\mathrm{c}}-F_{\mathrm{n}}}{k_{\mathrm{B}}T}\right),\quad \frac{1-f_{\mathrm{p}}}{f_{\mathrm{p}}} & =\exp\left(\frac{F_{\mathrm{p}}-E_{\mathrm{v}}}{k_{\mathrm{B}}T}\right) \\
	\label{eq:f_n_identity}
\end{align}

The matrix elements cancel in the quotient of the absorption and stimulated-emission
rates, so that it equals the ratio of the occupation factors
$\Phi^{\mathrm{abs}}/\Phi^{\mathrm{stim}}$. Since the $\delta$-function fixes $E_{\mathrm{c}}
-E_{\mathrm{v}}=\hbar\omega$, it follows that:
\begin{equation}
	\frac{R^{\mathrm{abs}}}{R^{\mathrm{stim}}}=\frac{(1-f_{\mathrm{n}})(1-f_{\mathrm{p}})}{f_{\mathrm{n}}f_{\mathrm{p}}}
	=\exp\left[\frac{\hbar\omega-(F_{\mathrm{n}}-F_{\mathrm{p}})}{k_{\mathrm{B}}T}\right
	] \label{eq:abs_stim_rate}
\end{equation}

In thermal equilibrium, $F_{\mathrm{n}}=F_{\mathrm{p}}$ and $R^{\mathrm{abs}}> R^{\mathrm{stim}}$:
absorption always dominates. However, wave amplification is present if
$R^{\mathrm{stim}}> R^{\mathrm{abs}}$, i.e., if $F_{\mathrm{n}}-F_{\mathrm{p}}>\hbar
\omega$. Laser operation therefore requires a quasi-Fermi level separation that
exceeds the photon energy, which can only be achieved by strong carrier injection.

\subsubsection{Stimulated vs Spontaneous Emission}

The quotient of the stimulated and spontaneous emission rates equals the photon
occupation number $n_{\mathrm{ph}}$ per mode. With the mode density $N$ and the spectral
photon density $S=N\,n_{\mathrm{ph}}$, it follows that:
\begin{equation}
	\frac{R^{\mathrm{stim}}}{R^{\mathrm{spon}}}=n_{\mathrm{ph}}(\hbar \omega)=\frac{S(\hbar\omega)}{N(\hbar\omega)}
	=\frac{\pi^{2}\hbar^{3}c^{3}}{n_{\mathrm{r}}^{3}}\,\frac{S(\hbar\omega)}{(\hbar\omega)^{2}}
	\label{eq:stim_spon_rate}
\end{equation}

Since the ratio $R^{\mathrm{stim}}/R^{\mathrm{spon}}$ grows linearly with $S$, a high
photon density, achieved by waveguiding and a Fabry-Perot Resonator, shifts the balance
towards stimulated emission.

\subsubsection{Thermal Equilibrium}

In thermal equilibrium, the quasi-Fermi levels coincide,
$F_{\mathrm{n}}=F_{\mathrm{p}}$, and the photon occupation number follows the Bose-Einstein
distribution. Using \cref{eq:abs_stim_rate} and \cref{eq:stim_spon_rate}, the rates
satisfy:
\begin{align}
	\frac{R^{\mathrm{abs}}-R^{\mathrm{spon}}}{R^{\mathrm{stim}}} & =\exp\left( \frac{\hbar \omega- (F_{\mathrm{n}}-F_{\mathrm{p}})}{k_{\mathrm{B}}T}\right) - \frac{1}{n_{\mathrm{ph}}}                \\
	                                                             & =\exp\left( \frac{\hbar \omega}{k_{\mathrm{B}}T}\right) - \left[\exp\left( \frac{\hbar \omega}{k_{\mathrm{B}}T}\right) - 1\right]=1
\end{align}

Hence, absorption and total emission are in balance
$R^{\mathrm{abs}}=R^{\mathrm{stim}}+R^{\mathrm{spon}}$.

\subsection{Spontaneous Emission Line Shape}
\label{sec:spontaneous_emission}

Following \cref{sec:line_shape}, the spontaneous emission rate can be expressed
using the equilibrium absorption coefficient (see \cref{sec:absorption_coefficient_and_gain}),
the Photon Density of States $N(\hbar \omega)$ and the occupation probabilities $f
_{\mathrm{n}}$, $f_{\mathrm{p}}$.
\begin{equation}
	R^{\mathrm{spon}}(\hbar \omega)=N(\hbar\omega) \frac{c}{n_{\mathrm{r}}}\alpha_{0}
	(\hbar \omega) f_{\mathrm{n}}f_{\mathrm{p}}
\end{equation}

The occupation factors are evaluated in the Boltzmann approximation:
\begin{align}
	f_{\mathrm{n}}(E) & =\frac{1}{1+\exp\left( \frac{E-F_{\mathrm{n}}}{k_{\mathrm{B}}T} \right)}\simeq \exp\left( \frac{F_{\mathrm{n}}-E}{k_{\mathrm{B}}T}\right) \\
	f_{\mathrm{p}}(E) & =\frac{1}{1+\exp\left( \frac{F_{\mathrm{p}}-E}{k_{\mathrm{B}}T} \right)}\simeq \exp\left( \frac{E-F_{\mathrm{p}}}{k_{\mathrm{B}}T}\right)
\end{align}

The factors $f_{\mathrm{n}}$ and $f_{\mathrm{p}}$ are evaluated at the band
energies $E_{\mathrm{c}}(k_{0})$ and $E_{\mathrm{v}}(k_{0})$ of the transition, the
$\delta$-function fixes their difference to $E_{\mathrm{c}}(k_{0})-E_{\mathrm{v}}
(k_{0})=\hbar\omega$. The individual band energies therefore drop out of the
product and only the quasi-Fermi-level splitting
$\Delta F=F_{\mathrm{n}}-F_{\mathrm{p}}$ and the photon energy remain:
\begin{align}
	f_{\mathrm{n}}f_{\mathrm{p}} & \simeq \exp \left( \frac{F_{\mathrm{n}}-F_{\mathrm{p}}-\left(E_{\mathrm{c}}(k_{0})-E_{\mathrm{v}}(k_{0})\right)}{k_{\mathrm{B}}T}\right)=\exp \left( \frac{\Delta F -\hbar \omega}{k_{\mathrm{B}}T}\right) \\
	                             & =\exp\left( \frac{\Delta F-E_{\mathrm{g}}}{k_{\mathrm{B}}T}\right) \exp\left( -\frac{\hbar \omega-E_{\mathrm{g}}}{k_{\mathrm{B}}T}\right)
\end{align}

Inserting this into the rate gives the analytical spontaneous emission spectrum
(for $\hbar\omega>E_{\mathrm{g}}$), where $\Delta F$ only enters the amplitude:
\begin{equation}
	R^{\mathrm{spon}}(\hbar \omega)= N(\hbar\omega)\frac{c}{n_{\mathrm{r}}}\,\alpha
	_{0}(\hbar \omega)\,\exp\left( \frac{\Delta F-E_{\mathrm{g}}}{k_{\mathrm{B}}T}\right
	) \exp\left( -\frac{\hbar \omega-E_{\mathrm{g}}}{k_{\mathrm{B}}T}\right)
\end{equation}

With $N\propto(\hbar\omega)^{2}$ and $\alpha_{0}\propto\sqrt{\hbar\omega-E_{\mathrm{g}}}
\,/\,\hbar\omega$, the spectral shape follows:
\begin{equation}
	R^{\mathrm{spon}}(\hbar \omega)\propto \hbar\omega\,\sqrt{\hbar\omega-E_{\mathrm{g}}}
	\,\exp\left( -\frac{\hbar \omega-E_{\mathrm{g}}}{k_{\mathrm{B}}T}\right)
\end{equation}

\subsubsection{Measurable Quantities}

The number of photons emitted per volume and time with photon energies $E=\hbar\omega$
between $E_{1}$ and $E_{2}$ is:
\begin{equation}
	N(E_{1}, E_{2}) = \int_{E_1}^{E_2}R^{\mathrm{spon}}(E) \, \mathrm{d}E
\end{equation}

Weighting each photon by its energy gives the spectral power density $p_{E}(E)=E\,
R^{\mathrm{spon}}(E)$ and the emitted optical power per volume:
\begin{equation}
	P(E_{1}, E_{2})=\int_{E_1}^{E_2}p_{E}(E) \, \mathrm{d}E = \int_{E_1}^{E_2}E \, R
	^{\mathrm{spon}}(E) \, \mathrm{d}E
\end{equation}

The extra factor $E$ turns the factor $\hbar\omega$ of the rate into $(\hbar\omega
)^{2}$ in the power spectrum:
\begin{equation}
	p_{E}(E)\propto (\hbar\omega)^{2}\sqrt{ \hbar \omega-E_{\mathrm{g}}}\exp\left( -
	\frac{\hbar \omega-E_{\mathrm{g}}}{k_{\mathrm{B}}T}\right)
\end{equation}

Spectrometers usually resolve the emission in wavelength rather than in energy.
With $E=hc/\lambda$ (vacuum wavelength $\lambda$), the integral transforms as:
\begin{align}
	P(\lambda_{1}, \lambda_{2})           & =\int_{E_1}^{E_2}p_{E}(E)\,\mathrm{d}E=\int_{\lambda_2}^{\lambda_1}\underbrace{p_{E}\!\left(\frac{hc}{\lambda}\right)\frac{hc}{\lambda^{2}}}_{p_{\lambda}(\lambda)}\,\mathrm{d}\lambda                          \\
	\Rightarrow\quad p_{\lambda}(\lambda) & =p_{E}\!\left(\frac{hc}{\lambda}\right)\frac{hc}{\lambda^{2}}\propto\frac{1}{\lambda^{4}}\sqrt{\frac{hc}{\lambda}-E_{\mathrm{g}}}\,\exp\left( -\frac{\frac{hc}{\lambda}-E_{\mathrm{g}}}{k_{\mathrm{B}}T}\right)
\end{align}

A real detector measures the power within a finite bandwidth $\Delta\lambda=\lambda
_{1}-\lambda_{2}$. If $p_{\lambda}(\lambda)$ varies only slightly within this
interval, the integral can be approximated by the midpoint rule, and the
spectral power density at the center wavelength
$\lambda_{\mathrm{c}}=(\lambda_{1}+\lambda_{2})/2$ is approximately the measured
power divided by the interval width:
\begin{equation}
	P(\lambda_{1},\lambda_{2})\approx p_{\lambda}(\lambda_{\mathrm{c}})\,\Delta\lambda
	\quad\Rightarrow\quad p_{\lambda}(\lambda_{\mathrm{c}})\approx\frac{P(\lambda_{1},\lambda_{2})}{\Delta\lambda}
\end{equation}